\chapter{Mathematical Preliminaries}
\label{Section:Preliminaries}

This chapter collects the tools used across the optimization and control
analyses. It presents standard material for which we make no claim of novelty.
This chapter was generated primarily using AI.

\section{Notation and online comparisons}\label{sec:prelim:notation}

We follow the notation of Section~\ref{sec:intro:notation}.
In this chapter, $Q\subseteq\Real^n$ is a convex feasible set and $T$ is
the horizon.

In online convex optimization, $z_t\in Q$ is selected before the convex
loss $\ell_t$ is revealed. Static regret on a fixed convex set $Q$
\cite{zinkevich2003online,hazan2016introduction} is
\begin{equation}\label{eq:intro:static-regret}
 \Reg_T(Q)\coloneqq\sum_{t=1}^T\ell_t(z_t)-\inf_{z\in Q}\sum_{t=1}^T\ell_t(z).
\end{equation}
On a nonempty, closed, bounded convex domain with uniformly bounded
subgradients, projected subgradient descent with suitable step sizes gives
$O(\sqrt T)$ regret.

In the storage-control chapter, $z$ specifies a fixed policy. To define
$\ell_t(z)$, we apply that policy from the first round using the same
arrivals and cost functions as the online controller. Here, $\ell_t(z)$ is the
resulting cost in round $t$. The online controller changes its policy
over time, so its current state also reflects earlier policies. Its
actual cost can therefore differ from $\ell_t(z_t)$. 


For a norm $\norm{\cdot}$, its dual is
\begin{equation}\label{eq:intro:dual-norm}
 \norm{g}_*\coloneqq\sup_{\norm{x}\le1}\inner{g}{x}.
\end{equation}
In particular, $\ell_1$ and $\ell_\infty$ are dual, and $\ell_2$ is self-dual.

\begin{Thm:Lemma}[H\"older and Young inequalities; cf.~{\cite[Example~5.11]{orabona_modern_2023}}]\label{lem:holder}
For any dual pair of norms and any $\gamma>0$,
\begin{equation}\label{eq:intro:holder-young}
 |\inner{g}{v}|\le\norm{g}_*\norm{v},\qquad
 \inner{g}{v}-\frac\gamma2\norm{v}^2
 \le\frac{\norm{g}_*^2}{2\gamma}.
\end{equation}
For the norm induced by the inner product, the first inequality is Cauchy--Schwarz.
For $a,b\ge0$ and conjugate exponents $p,q>1$, meaning $1/p+1/q=1$,
Young's inequality also gives
\begin{equation}\label{eq:intro:young-conjugate}
 ab\le\frac{a^p}{p}+\frac{b^q}{q}.
\end{equation}
\end{Thm:Lemma}
\begin{proof}
For $v\ne0$, apply~\eqref{eq:intro:dual-norm} to both
$v/\norm{v}$ and $-v/\norm{v}$. The case $v=0$ is immediate.
With $r\coloneqq\norm{v}$, this gives
\begin{align*}
 \inner{g}{v}-\frac\gamma2\norm{v}^2
 &\le \norm{g}_*r-\frac\gamma2r^2\\
 &=-\frac\gamma2\left(r-\frac{\norm{g}_*}{\gamma}\right)^2
   +\frac{\norm{g}_*^2}{2\gamma}
 \le \frac{\norm{g}_*^2}{2\gamma}.
\end{align*}
For~\eqref{eq:intro:young-conjugate}, the case $b=0$ is immediate.
If $b>0$, the derivative of $br-r^p/p$ is $b-r^{p-1}$, so its
maximum over $r\ge0$ occurs at $r=b^{1/(p-1)}$. Since $q=p/(p-1)$,
\[
 ab-\frac{a^p}{p}
 \le\max_{r\ge0}\left\{br-\frac{r^p}{p}\right\}
 =b^q-\frac{b^q}{p}=\frac{b^q}{q}.
 \qedhere
\]
\end{proof}

\section{Convexity, subgradients, and affine parameterizations}
\label{sec:intro:convexity}

A set $Q$ is convex if it contains every segment between its points.
A function $f:Q\to\Real$ is convex if
\begin{equation}\label{eq:intro:convexity}
 f(\theta x+(1-\theta)y)\le\theta f(x)+(1-\theta)f(y),
 \qquad x,y\in Q,\quad 0\le\theta\le1.
\end{equation}
A subgradient $g\in\partial f(x)$ satisfies
\begin{equation}\label{eq:intro:subgradient}
 f(y)\ge f(x)+\inner{g}{y-x},\qquad y\in Q.
\end{equation}
When $f$ is differentiable at $x$, we use $\nabla f(x)$ as a subgradient.

\begin{Thm:Lemma}[Convex comparison and affine composition {\cite{boyd_vandenberghe_2004,orabona_modern_2023}}]
\label{lem:intro:convex-comparison}
Let $f$ be convex on $Q$.
For $x_i\in Q$, $\theta_i\ge0$, and $\sum_i\theta_i=1$,
\begin{equation}\label{eq:intro:jensen}
 f\Bigl(\sum_i\theta_i x_i\Bigr)\le\sum_i\theta_i f(x_i).
\end{equation}
For affine composition, let $P\subseteq\Real^p$ and $Q\subseteq\Real^q$
be convex, with $A\in\Real^{q\times p}$ and $b\in\Real^q$.
If $F(z)\coloneqq Az+b$ maps $P$ into $Q$, then $\ell(z)\coloneqq f(F(z))$ is convex and
\begin{equation}\label{eq:intro:affine-subgradient}
 g\in\partial f(F(z))\quad\Longrightarrow\quad
 A^\top g\in\partial\ell(z).
\end{equation}
\end{Thm:Lemma}
\begin{proof}
For~\eqref{eq:intro:jensen}, discard zero weights and induct on the number
of remaining weights. If there is more than one, choose a weight
$0<\theta_j<1$, normalize the others by $1-\theta_j$, and apply the
induction hypothesis to their convex combination, followed by
\eqref{eq:intro:convexity} with weights $\theta_j$ and $1-\theta_j$.
For $z,z'\in P$ and $0\le\lambda\le1$,
\[
 \ell(\lambda z+(1-\lambda)z')
 =f\bigl(\lambda F(z)+(1-\lambda)F(z')\bigr)
 \overset{\mathrm{(A)}}{\le}\lambda\ell(z)+(1-\lambda)\ell(z'),
\]
where (A) follows from~\eqref{eq:intro:convexity}, since
$F(z),F(z')\in Q$.
For every $z'\in P$,
\[
 \ell(z')-\ell(z)
 =f(F(z'))-f(F(z))
 \overset{\mathrm{(A)}}{\ge}\inner{g}{F(z')-F(z)}
 =\inner{A^\top g}{z'-z}.
\]
Here (A) follows from~\eqref{eq:intro:subgradient},
proving~\eqref{eq:intro:affine-subgradient}.
\end{proof}

\begin{Thm:Lemma}[Jensen's inequality for fractional powers]
\label{lem:intro:jensen-expectation}
For a nonnegative random variable $Z$ with $\E[Z]<\infty$ and $0<\theta<1$,
\begin{equation}\label{eq:intro:jensen-power}
 \E[Z^\theta]\le(\E[Z])^\theta.
\end{equation}
\end{Thm:Lemma}
\begin{proof}
Put $\mu\coloneqq\E[Z]$.
If $\mu=0$, then $Z=0$ almost surely and the conclusion is immediate.
If $\mu>0$, concavity of $z\mapsto z^\theta$ on $[0,\infty)$ gives
\[
 z^\theta\le\mu^\theta+\theta\mu^{\theta-1}(z-\mu),\qquad z\ge0.
\]
Taking expectations with $z=Z$ proves~\eqref{eq:intro:jensen-power}.
\end{proof}

\begin{Thm:Lemma}[Bounded subgradient selections and Lipschitz continuity; cf.~{\cite[Theorem~2.24]{orabona_modern_2023}}]
\label{lem:intro:lipschitz}
Suppose that at every $x\in Q$ a subgradient $g_x\in\partial f(x)$ can be
selected with $\norm{g_x}_*\le L$. Then
\begin{equation}\label{eq:intro:lipschitz}
 |f(x)-f(y)|\le L\norm{x-y},\qquad x,y\in Q.
\end{equation}
\end{Thm:Lemma}
\begin{proof}
For $x,y\in Q$,
\[
 f(x)-f(y)
 \le\inner{g_x}{x-y}
 \le\norm{g_x}_*\norm{x-y}
 \le L\norm{x-y}.
\]
Interchanging $x$ and $y$ proves~\eqref{eq:intro:lipschitz}.
\end{proof}

Inequality~\eqref{eq:intro:lipschitz} defines $L$-Lipschitz continuity.

In the mirror-descent updates below, strong convexity of the regularizer
bounds the distance between successive iterates.
See~\eqref{eq:intro:md-movement}.

\begin{Thm:Definition}[Strong convexity {\cite[Section~4.1.1]{orabona_modern_2023}}]\label{def:strong-convex}
A function $f:Q\to\Real$ is $\sigma$-strongly convex with respect to
$\norm{\cdot}$ if, for $x,y\in Q$ and $0\le\theta\le1$,
\begin{equation}\label{eq:intro:strong-convexity}
 f(\theta x+(1-\theta)y)
 \le\theta f(x)+(1-\theta)f(y)
 -\frac\sigma2\theta(1-\theta)\norm{x-y}^2.
\end{equation}
\end{Thm:Definition}

If $f$ is $\sigma$-strongly convex on $Q$, then for $x\in Q$ and
$g\in\partial f(x)$,
\begin{equation}\label{eq:intro:strong-support}
 f(y)\ge f(x)+\inner{g}{y-x}
 +\frac\sigma2\norm{y-x}^2,
 \qquad y\in Q.
\end{equation}
\begin{proof}
For $x_\theta\coloneqq(1-\theta)x+\theta y$ and $0<\theta<1$,
\begin{align*}
 f(x)+\theta\inner{g}{y-x}
 &\overset{\mathrm{(A)}}{\le} f(x_\theta)\\
 &\overset{\mathrm{(B)}}{\le}(1-\theta)f(x)+\theta f(y)
   -\frac\sigma2\theta(1-\theta)\norm{x-y}^2.
\end{align*}
Here (A) uses~\eqref{eq:intro:subgradient} with comparison point
$x_\theta\in Q$. Step (B) uses
\eqref{eq:intro:strong-convexity} with weights $1-\theta$ and $\theta$.
Rearranging, dividing by $\theta>0$, and letting
$\theta\downarrow0$ proves
\eqref{eq:intro:strong-support}.
\end{proof}

\section{Bregman divergence and the three-point inequality}
\label{sec:intro:bregman}

A Bregman divergence serves as a proximity term in mirror
descent \cite{shalev2007primal,abernethy2009competing,hazan2010extracting}.
The three-point inequality below converts the minimization that defines an
update into a bound relative to any feasible comparison point
\cite{prox-siam93,SGD-robust,tseng-accelerated,wei2020online}.

Let $C$ be a convex domain and $\Phi:C\to\Real$ a convex function,
differentiable on $\operatorname{ri}(C)$, the interior relative
to its affine hull. Derivatives are taken in that affine hull. The
constraint set $Q\subseteq C$ can have smaller dimension. Define
$Q_{\mathrm{ref}}\coloneqq Q\cap \operatorname{ri}(C)$. For $x\in Q$, $y\in Q_{\mathrm{ref}}$, the Bregman
divergence~\cite{bregman1967relaxation} is
\begin{equation}\label{eq:intro:bregman}
 D_\Phi(x,y)\coloneqq\Phi(x)-\Phi(y)-\inner{\nabla\Phi(y)}{x-y}.
\end{equation}
The restriction on the second argument ensures that $\nabla\Phi(y)$ is
defined. The first argument may lie on the boundary, because the formula
requires only the value $\Phi(x)$ there. For entropy, this allows comparison
points with zero coordinates while requiring positive coordinates at the
reference point.
Applying~\eqref{eq:intro:subgradient} to $\Phi$ at $y$ gives
$D_\Phi(x,y)\ge0$. If $\Phi$ is $\sigma$-strongly convex on $Q$, then
\eqref{eq:intro:strong-support} with base point $y$, comparison point $x$,
and subgradient $\nabla\Phi(y)$ gives
\begin{equation}\label{eq:intro:bregman-curvature}
 D_\Phi(x,y)\ge\frac\sigma2\norm{x-y}^2,
 \qquad x\in Q,\quad y\in Q_{\mathrm{ref}}.
\end{equation}
Expanding each divergence using~\eqref{eq:intro:bregman} gives the
three-point identity
\begin{equation}\label{eq:intro:three-point}
 \inner{\nabla\Phi(x)-\nabla\Phi(y)}{z-x}
 =D_\Phi(z,y)-D_\Phi(z,x)-D_\Phi(x,y),
\end{equation}
for $x,y\in Q_{\mathrm{ref}}$ and $z\in Q$.

\begin{Thm:Lemma}[Bregman three-point inequality; cf.~{\cite[Lemma~2.1]{wei2020online}}]\label{lem:pushback}
Let $Q\subseteq C$ be convex, let $f:Q\to\Real$ be convex, and fix
$\alpha>0$ and $y\in Q_{\mathrm{ref}}$. Suppose
\begin{equation}\label{eq:bregman-min}
 \hat x\in\argmin_{x\in Q}\{f(x)+\alpha D_\Phi(x,y)\},
 \qquad \hat x\in Q_{\mathrm{ref}}.
\end{equation}
Then, for all $z\in Q$,
\begin{equation}\label{eq:pushback}
 f(\hat x)+\alpha D_\Phi(\hat x,y)
 \le f(z)+\alpha D_\Phi(z,y)-\alpha D_\Phi(z,\hat x).
\end{equation}
\end{Thm:Lemma}
\begin{proof}
For $x_\theta\coloneqq(1-\theta)\hat x+\theta z$ and $0<\theta<1$,
convexity of $Q$ gives $x_\theta\in Q$, so
\begin{align*}
 f(\hat x)+\alpha D_\Phi(\hat x,y)
 &\overset{\mathrm{(A)}}{\le} f(x_\theta)+\alpha D_\Phi(x_\theta,y)\\
 &\overset{\mathrm{(B)}}{\le}(1-\theta)f(\hat x)+\theta f(z)
   +\alpha D_\Phi(x_\theta,y).
\end{align*}
Here (A) is optimality in~\eqref{eq:bregman-min}, and (B) is
\eqref{eq:intro:convexity} for $f$ with weights $1-\theta$ and $\theta$.
Rearranging and dividing by $\theta>0$ gives
\[
 f(\hat x)-f(z)
 \le\alpha\frac{D_\Phi(x_\theta,y)-D_\Phi(\hat x,y)}{\theta}.
\]
By~\eqref{eq:intro:bregman}, the difference quotient on the right tends to
$\inner{\nabla\Phi(\hat x)-\nabla\Phi(y)}{z-\hat x}$ as
$\theta\downarrow0$. Thus
\[
 f(\hat x)-f(z)
 \le\alpha\inner{\nabla\Phi(\hat x)-\nabla\Phi(y)}{z-\hat x}
 =\alpha\{D_\Phi(z,y)-D_\Phi(z,\hat x)-D_\Phi(\hat x,y)\}.
 \qedhere
\]
\end{proof}

\section{Quadratic and entropy geometries}\label{sec:intro:geometry}

We use quadratic regularizers for Euclidean updates,
and entropy regularizers for allocation fractions on a simplex.

For the Euclidean norm, set
\begin{equation}\label{eq:intro:quadratic-geometry}
 \Phi(x)\coloneqq\tfrac12\norm{x}^2,\qquad
 D_\Phi(x,y)=\tfrac12\norm{x-y}^2.
\end{equation}
The function $\Phi$ is $1$-strongly convex in the Euclidean norm, whose
dual norm is also Euclidean.

For entropy, define the probability and subprobability simplices
\begin{equation}\label{eq:intro:simplices}
 \Delta_d\coloneqq\{p\ge0:\mathbf1^\top p=1\},\qquad
 S_d\coloneqq\{p\ge0:\mathbf1^\top p\le1\},
\end{equation}
and the unnormalized negative entropy
\begin{equation}\label{eq:intro:negentropy}
 \Psi(p)\coloneqq\sum_{i=1}^d(p_i\log p_i-p_i),\qquad 0\log0\coloneqq0.
\end{equation}
Its Bregman divergence is the generalized Kullback--Leibler (KL) divergence:
for $p\ge0$ and $q>0$,
\begin{equation}\label{eq:intro:generalized-kl}
 D_\Psi(p,q)=\sum_i\left(p_i\log\frac{p_i}{q_i}-p_i+q_i\right).
\end{equation}
When $p,q\in\Delta_d$ with $q>0$, the linear terms cancel, giving the
ordinary KL divergence between probability vectors.

\begin{Thm:Lemma}[Bounds on entropy divergence; cf.~{\cite[Proposition~5.1]{BECK2003167}}]
\label{lem:intro:entropy}
For $p\in S_d$ and $q\in S_d\cap(0,\infty)^d$,
\begin{equation}\label{eq:intro:entropy-curvature}
 D_\Psi(p,q)\ge\tfrac12\norm{p-q}_1^2.
\end{equation}
The uniform initial points satisfy
\begin{equation}\label{eq:intro:entropy-initial}
 \sup_{p\in\Delta_d}D_\Psi(p,\mathbf1/d)\le\log d,
 \qquad
 \sup_{p\in S_d}D_\Psi(p,\mathbf1/(d+1))\le\log(d+1).
\end{equation}
\end{Thm:Lemma}
\begin{proof}
For positive $p,q\in S_d$, put $v\coloneqq p-q$ and $q_s\coloneqq q+sv$ for
$0\le s\le1$. Then $q_s>0$, $q_s\in S_d$, and differentiating
\eqref{eq:intro:negentropy} gives
$\nabla^2\Psi(q_s)=\operatorname{diag}(1/q_{s,i})$. Hence
\begin{equation}\label{eq:intro:entropy-hessian-bound}
  \norm{v}_1^2
 \overset{\mathrm{(A)}}{\le}
 \left(\sum_i\frac{v_i^2}{q_{s,i}}\right)\sum_iq_{s,i}
 \overset{\mathrm{(B)}}{\le} v^\top\nabla^2\Psi(q_s)v.
\end{equation}
Step (A) is Cauchy--Schwarz from Lemma~\ref{lem:holder}, applied to
the vectors with components $|v_i|/\sqrt{q_{s,i}}$ and
$\sqrt{q_{s,i}}$.
Step (B) uses $\sum_iq_{s,i}\le1$ from~\eqref{eq:intro:simplices}.
Taylor's formula with integral remainder yields
\[
 D_\Psi(p,q)
 =\int_0^1(1-s)v^\top\nabla^2\Psi(q_s)v\,\mathrm{d}s
 \overset{\mathrm{(A)}}{\ge}\norm{v}_1^2\int_0^1(1-s)\,\mathrm{d}s
 =\frac12\norm{p-q}_1^2.
\]
Here (A) integrates~\eqref{eq:intro:entropy-hessian-bound} against the
nonnegative weight $1-s$.
The same bound holds when $p$ has zero coordinates by continuity, using
$0\log0=0$.
For the initial bounds, $D_\Psi(\cdot,q)$ is convex for fixed $q>0$.
Writing $e_i$ for the coordinate vectors, every point of $\Delta_d$ is a
convex combination of the $e_i$, and every point of $S_d$ is a convex
combination of $0,e_1,\ldots,e_d$. Thus~\eqref{eq:intro:jensen} reduces
the bounds to the vertex values from~\eqref{eq:intro:generalized-kl}:
\[
\begin{aligned}
 D_\Psi(e_i,\mathbf1/d)&=\log d,\\
 D_\Psi(0,\mathbf1/(d+1))&=\frac d{d+1},\\
 D_\Psi(e_i,\mathbf1/(d+1))&=\log(d+1)-\frac1{d+1}.
\end{aligned}
\]
Since $d/(d+1)\le\int_1^{d+1}t^{-1}\,\mathrm{d}t=\log(d+1)$,
these values prove~\eqref{eq:intro:entropy-initial}, including $d=1$.
\end{proof}

For the uniform initial points in~\eqref{eq:intro:entropy-initial},
the term involving the initial Bregman divergence in the mirror-descent bound
\eqref{eq:intro:md-regret} is at most $\log d/\eta$ on $\Delta_d$ and
$\log(d+1)/\eta$ on $S_d$, where $\eta>0$ is the step size.
We measure subgradients in the $\ell_\infty$ norm, dual to $\ell_1$.

\begin{Thm:Lemma}[Entropy mirror-descent updates; cf.~{\cite[Section~5]{BECK2003167}}]\label{lem:intro:entropy-update}
\label{ch4:lem:omd-closed-form}
Let $q$ be positive in $Q=\Delta_d$ or $Q=S_d$, let $\eta>0$, and
set $a_i\coloneqq q_i e^{-\eta g_i}$. The unique minimizer of
$\eta\inner{g}{p}+D_\Psi(p,q)$ over $p\in Q$ is
\begin{equation}\label{eq:intro:entropy-update}
 p^+\coloneqq
 \begin{cases}
 a/(\sum_i a_i),&Q=\Delta_d,\\
 a/\max\{1,\sum_i a_i\},&Q=S_d.
 \end{cases}
\end{equation}
Both updates preserve positivity.
\end{Thm:Lemma}
\begin{proof}
Let $A\coloneqq\sum_i a_i>0$ and let $p^+$ be the candidate in
\eqref{eq:intro:entropy-update}. Write $p^+=a/Z$, where $Z\coloneqq A$ on
$\Delta_d$ and $Z\coloneqq\max\{1,A\}$ on $S_d$, and put $\lambda\coloneqq\log Z$.
Both candidates are positive and feasible. Since $\log a_i=\log q_i-\eta g_i$,
\eqref{eq:intro:generalized-kl} gives
\[
 \eta\inner{g}{p}+D_\Psi(p,q)
 =D_\Psi(p,a)+\sum_i(q_i-a_i).
\]
The last sum is independent of $p$. For every feasible $p$, expansion gives
\[
 D_\Psi(p,a)-D_\Psi(p^+,a)
 =D_\Psi(p,p^+)+\lambda\bigl(\mathbf1^\top p^+-\mathbf1^\top p\bigr).
\]
This identity also holds when $p$ has zero coordinates, with $0\log0=0$.
On $\Delta_d$, the last term vanishes because both masses equal one.
On $S_d$, either $A\le1$ and $\lambda=0$, or $A>1$, $\lambda>0$, and
$\mathbf1^\top p^+=1\ge\mathbf1^\top p$. The difference is therefore at
least $D_\Psi(p,p^+)\ge\tfrac12\norm{p-p^+}_1^2$ by
\eqref{eq:intro:entropy-curvature}. Hence $p^+$ is the unique minimizer.
\end{proof}

\section{Mirror descent}
\label{sec:intro:mirror-descent}

The movement bound below is used in the storage-control chapter to account for
changes in the controller's policy.

\begin{Thm:Lemma}[Mirror descent: comparison, regret, and movement {\cite{BECK2003167,bubeck2015convex}}]
\label{lem:intro:mirror-descent}
Suppose $\Phi$ is $\sigma$-strongly convex on a convex set $Q$, where
$\sigma>0$, and $x,x^+\in Q_{\mathrm{ref}}$ satisfy
\begin{equation}\label{eq:intro:md-update}
 x^+\in\argmin_{v\in Q}\{\eta\inner{g}{v}+D_\Phi(v,x)\},
 \qquad\eta>0.
\end{equation}
For every $u\in Q$,
\begin{align}
 \inner{g}{x-u}
 &\le\frac{D_\Phi(u,x)-D_\Phi(u,x^+)}{\eta}
       +\frac\eta{2\sigma}\norm{g}_*^2,
       \label{eq:intro:md-one-step}\\
 \norm{x^+-x}&\le\frac\eta\sigma\norm{g}_*.
       \label{eq:intro:md-movement}
\end{align}
For a sequence of these updates with a common step size and
$g_t\in\partial\ell_t(x_t)$, convexity of $\ell_t$ gives
\begin{equation}\label{eq:intro:md-regret}
 \sum_{t=1}^T\bigl(\ell_t(x_t)-\ell_t(u)\bigr)
 \le\frac{D_\Phi(u,x_1)}\eta
       +\frac\eta{2\sigma}\sum_{t=1}^T\norm{g_t}_*^2.
\end{equation}
The same bound holds for the sum of $\inner{g_t}{x_t-u}$.
\end{Thm:Lemma}
\begin{proof}
Apply Lemma~\ref{lem:pushback} with $f(v)\coloneqq\eta\inner{g}{v}$,
$\alpha=1$, $y=x$, $\hat x=x^+$, and comparison point $z=u$.
Both reference points are admissible by $x,x^+\in Q_{\mathrm{ref}}$, so
\begin{equation}\label{eq:intro:md-pushback}
 \eta\inner{g}{x^+-u}
 \le D_\Phi(u,x)-D_\Phi(u,x^+)-D_\Phi(x^+,x).
\end{equation}
Adding $\eta\inner{g}{x-x^+}$ and applying H\"older's inequality,
the curvature bound~\eqref{eq:intro:bregman-curvature}, and Young's
inequality~\eqref{eq:intro:holder-young} with $\gamma=\sigma$ gives
\begin{align*}
 \eta\inner{g}{x-u}
 &\le D_\Phi(u,x)-D_\Phi(u,x^+)
   +\eta\norm{g}_*\norm{x-x^+}
   -\frac\sigma2\norm{x-x^+}^2\\
 &\le D_\Phi(u,x)-D_\Phi(u,x^+)
   +\frac{\eta^2}{2\sigma}\norm{g}_*^2.
\end{align*}
Dividing by $\eta>0$ proves~\eqref{eq:intro:md-one-step}.
For movement, set $u=x$ in~\eqref{eq:intro:md-pushback} and apply
\eqref{eq:intro:bregman-curvature} in both directions:
\[
  \sigma\norm{x^+-x}^2
 \le D_\Phi(x,x^+)+D_\Phi(x^+,x)
 \le\eta\inner{g}{x-x^+}
 \le\eta\norm{g}_*\norm{x^+-x}.
\]
The last inequality is H\"older's. If $x^+=x$, the movement bound is
immediate. Otherwise, divide by $\sigma\norm{x^+-x}>0$ to obtain
\eqref{eq:intro:md-movement}.

For the sequence of updates, convexity and summation of
\eqref{eq:intro:md-one-step} with a common step size give
\begin{align*}
 \sum_{t=1}^T\bigl(\ell_t(x_t)-\ell_t(u)\bigr)
 &\le\sum_{t=1}^T\inner{g_t}{x_t-u}\\
 &\le\frac{D_\Phi(u,x_1)-D_\Phi(u,x_{T+1})}{\eta}
   +\frac\eta{2\sigma}\sum_{t=1}^T\norm{g_t}_*^2\\
 &\le\frac{D_\Phi(u,x_1)}\eta
   +\frac\eta{2\sigma}\sum_{t=1}^T\norm{g_t}_*^2.
\end{align*}
This proves~\eqref{eq:intro:md-regret} and
the same bound for the sum of inner products.
\end{proof}

\section{Projection and positive-part bounds}
\label{sec:intro:projection}

The projection condition is used to derive the domain-separation update.
The positive-part bounds below control clipped multipliers and storage
trajectories.

\begin{Thm:Lemma}[Euclidean projection {\cite{bubeck2015convex}}]\label{lem:intro:projection}
Let $Q\subseteq\Real^n$ be nonempty, closed, and convex, and define
$\Pi_Q(v)\coloneqq\argmin_{q\in Q}\norm{q-v}^2$. For $v\in\Real^n$ and $p\in Q$,
\begin{equation}\label{eq:intro:projection-optimality}
 p=\Pi_Q(v)
 \quad\Longleftrightarrow\quad
 \inner{p-v}{z-p}\ge0\qquad\text{for every }z\in Q.
\end{equation}
\end{Thm:Lemma}
\begin{proof}
If $p=\Pi_Q(v)$, the right directional derivative of
$\norm{p+\theta(z-p)-v}^2$ at $\theta=0$ is nonnegative for every
$z\in Q$, giving $2\inner{p-v}{z-p}\ge0$.
Conversely, the stated inequality gives
\[
 \norm{z-v}^2-\norm{p-v}^2
 =2\inner{p-v}{z-p}+\norm{z-p}^2
 \ge\norm{z-p}^2\ge0,
 \qquad z\in Q.
\]
Thus $p$ is the unique minimizer.
\end{proof}

% Keep the final lemma with its proof at the end of the chapter.
\begin{samepage}
\begin{Thm:Lemma}[Positive-part bounds]
\label{lem:intro:positive-part}
For $a,b\in\Real$,
\begin{equation}\label{eq:intro:positive-part-nonexpansive}
 |[a]_+-[b]_+|\le|a-b|.
\end{equation}
For $v\in\Real^n$, with the Euclidean norm and the positive part taken
componentwise,
\begin{equation}\label{eq:intro:positive-part-norm}
 \norm{[v]_+}\le\norm{v}.
\end{equation}
\end{Thm:Lemma}
\begin{proof}
By symmetry, assume $a\ge b$. Then
\[
 [a]_+-[b]_+
 =\begin{cases}
  0,&a\le0,\\
  a-b,&b\ge0,\\
  a,&b<0<a,
 \end{cases}
 \qquad 0\le[a]_+-[b]_+\le a-b.
\]
This proves~\eqref{eq:intro:positive-part-nonexpansive}. Applying it
with $b=0$ in each coordinate and summing squares gives
$\norm{[v]_+}^2\le\norm{v}^2$, hence
\eqref{eq:intro:positive-part-norm}.
\end{proof}
\end{samepage}
